\newpage

\section{Laborversuch}
Hier kommt alles zum dritten Laborversuch

\section*{Kurzzusammenfassung} 
        In diesem Laborversuch wurde eine astabile Multivibratorschaltung 
        mit einem LM358 Operationsvertärker aufgebaut. 
        Diese Schaltung erzeugt ein Rechtecksignal ohne ein externes Eingangssignal.
        Anschließend wurde die Bandbreite des LM358 bei einer Verstärkung von 1
        und die Bandbreite bei einer Verstärkung von 10 gemessen. Zum Schluss wurde die 
        Input-Offset-Voltage des OPV gemessen.
\section*{Durchführung} 
        \subsection*{1}
        Zu Beginn wurde die Schaltung des astabilen Multivibrator aufgebaut. 
        Dafür wurde eine negative Spannung von \qty{5}{\volt} benötigt. Diese wurde mithilfe 
        von einer 2-Kanal-Gleichsspannungsversorgung, welche galvanisch getrennte Kanäle besitzt, 
        erzeugt. Dabei wurde der positive Anschluss des Kanal 1 
        mit dem negativen Anschluss des Kanal 2 verbunden. Der positive Anschluss des Kanal 2 war somit \qty{5}{\volt}  
        und der negative Anschluss des Kanal 1 war somit \qty{-5}{\volt}. 
        Als Ground wurde der Ausgang der negative Ausgang des Kanal 2 verwendet. 
        Anschließend wurde die Schaltung aufgebaut und mit der Gleichspannungsquelle betrieben. 
        Mit dem Oszilloskop wurde die Spannung über den Kondensator gemessen und die Spannung am Ausgang des Operationsvertärker, siehe \ref{fig:7}.
\begin{figure}

        
        \label{fig:7}
\end{figure}
        Die Schaltung besteht aus zwei Teilen der Teil des OPV der in Mitkopplung betrieben wird
        ist ein Schmitt-Trigger. Dieser ermöglicht es das die Schaltung ab einer bestimmten Spannungsschwelle schaltet.
        Diese Spannungsschwelle berechnet sich mit folgender Formel:
        \begin{equation}
               $ r_1 = $
        \end{equation}



\section*{Diskussion}
